% appx-refs-manifest.tex — Appendix E: source-document manifest.

\chapter{Source-Document Manifest}
\label{ch:appx-refs-manifest}

\noindent\emph{The PDFs of the manuscript's primary literature
sources, with the chapters that cite each.}

This appendix is the curated source-document manifest: the
primary-literature PDFs the manuscript's analytical inputs draw
from, and the chapters where each is used.  The full BibTeX
bibliography is at the back of the monograph; this manifest is
the reader's quick orientation.

\section{Manifest}
\label{sec:refs-manifest}

\begin{table}[h]
\centering
\renewcommand{\arraystretch}{1.3}
\begin{tabular}{p{6.5cm}@{\hspace{1em}}p{5cm}@{\hspace{0.5em}}p{2.5cm}}
\toprule
Source                                                            & Cited in                                            & BibTeX key                  \\
\midrule
\textbf{Clay statement}
                                                                  &                                                     &                             \\
\quad Jaffe \& Witten,
``Quantum Yang--Mills Theory'' (2006)                              & \Cref{ch:introduction},
                                                                    \Cref{ch:mass-gap-theorem}                        &
\codepath{JaffeWitten2006}     \\
\midrule
\textbf{Axiomatic frameworks}
                                                                  &                                                     &                             \\
\quad Streater \& Wightman,
``PCT, Spin and Statistics, and All That''                         & \Cref{ch:introduction},
                                                                    \Cref{ch:os-axioms}                                &
\codepath{StreaterWightman2000}\\
\quad Osterwalder \& Schrader,
``Axioms for Euclidean Green's functions''                          & \Cref{ch:os-axioms},
                                                                    \Cref{ch:os-reconstruction}                        &
\codepath{OsterwalderSchrader1973}\\
\quad Glimm \& Jaffe,
``Quantum Physics: A Functional Integral Point of View''            & throughout (Parts III--V)                          &
\codepath{GlimmJaffe1987}      \\
\midrule
\textbf{Bałaban RG series}
                                                                  &                                                     &                             \\
\quad Bałaban, ``Renormalization group approach to lattice
gauge field theories. I'' (CMP 109, 1987)                          & \Cref{ch:small-field-kp},
                                                                    \Cref{ch:buffer-shell}                             &
\codepath{Balaban1987RG-I}     \\
\quad Bałaban, ``Renormalization group approach to lattice
gauge field theories. II'' (CMP 116, 1988)                         & \Cref{ch:small-field-kp},
                                                                    \Cref{ch:buffer-shell}                             &
\codepath{Balaban1988RG-II}    \\
\quad Bałaban, ``Convergent renormalization expansions for
lattice gauge theories'' (CMP 119, 1988)                           & \Cref{ch:small-field-kp},
                                                                    \Cref{ch:sigma-hat-trajectory}                     &
\codepath{Balaban1988convergent}\\
\midrule
\textbf{Cluster-expansion foundations}
                                                                  &                                                     &                             \\
\quad Brydges \& Federbush,
``A new form of the Mayer expansion'' (J. Math. Phys. 19, 1978)     & \Cref{ch:small-field-kp},
                                                                    \Cref{ch:large-field-kp}                           &
\codepath{BrydgesFederbush1978}\\
\quad Brydges, Fröhlich \& Seiler,
``Diamagnetic and critical properties of Higgs lattice gauge
theories'' (IHES, 1978)                                            & \Cref{ch:chaining-estimate}                       &
\codepath{BrydgesFroehlichSeiler1978IHES}\\
\quad Cammarota, ``Decay of correlations for infinite-range
interactions'' (CMP 85, 1982)                                      & \Cref{ch:small-field-kp}                          &
\codepath{Cammarota1982}       \\
\quad Seiler, ``Gauge theories as a problem of constructive QFT''
(LNP 159, 1982)                                                    & \Cref{ch:large-field-kp}                          &
\codepath{Seiler1982}          \\
\midrule
\textbf{Lattice gauge spectroscopy}
                                                                  &                                                     &                             \\
\quad Teper, ``$\SUN$ gauge theories in $(2+1)$ dimensions: in
search of clear results for theorists'' (hep-th/9812187, 1998)      & \Cref{ch:mass-gap-discussion}                     &
\codepath{Teper1998Glueballs}  \\
\quad Lucini, Teper \& Wenger,
``Glueballs and $k$-strings in $\SUN$ gauge theories''
(JHEP 2004)                                                        & \Cref{ch:mass-gap-discussion}                     &
\codepath{LuciniTeperWenger2004}\\
\midrule
\textbf{Heat-kernel and character expansions}
                                                                  &                                                     &                             \\
\quad Drouffe \& Itzykson, ``Lattice gauge fields''
(Phys. Rep. 38, 1978)                                              & \Cref{ch:heat-kernel},
                                                                    \Cref{ch:large-field-kp}                           &
\codepath{DrouffeItzykson1978} \\
\quad Itzykson \& Drouffe, \emph{Statistical Field Theory}
(Cambridge, 1989)                                                  & \Cref{ch:large-field-kp}                          &
\codepath{ItzyksonDrouffe1989} \\
\midrule
\textbf{Constructive analysis}
                                                                  &                                                     &                             \\
\quad Jörgensen, ``Analytic continuation of local representations
of Lie groups'' (1986)                                              & \Cref{ch:os-reconstruction}                       &
\codepath{Jorgensen1986AnalyticContinuation} \\
\bottomrule
\end{tabular}
\caption{Source-document manifest.  Every entry has a
corresponding PDF in \codepath{docs/pdf\_references/} for
project-internal audit; the full bibliographic detail is in the
manuscript's \texttt{biblatex} bibliography at the back.}
\label{tab:refs-manifest}
\end{table}

\section{Note on the audit trail}
\label{sec:refs-audit-trail}

\marginkey{Direct-PDF audit}%

Of the sources cited, several have been the subject of explicit
direct-PDF audits in the project's notes, with the audit findings
folded back into the constant provenance:

\begin{itemize}
  \item \codepath{Balaban1987RG-I} and \codepath{Balaban1988RG-II}
    were audited (2026-05-25) for the spectral-gap coefficient
    $\clambda$.  The audit confirmed Bałaban's framework
    establishes the \emph{existence} of $\clambda$ satisfying a
    system of upper-bound constraints, but does not pin a numerical
    value.  The manuscript's identification
    $\clambda = \Casimir(\tfrac12) = \tfrac34$ comes from the
    SU(2)-Casimir-gap argument in \Cref{ch:small-field-kp},
    Theorem~\ref{thm:c-lambda-casimir}, which is project-derived
    rather than directly extracted.
  \item \codepath{BrydgesFroehlichSeiler1978IHES} was audited
    (2026-05-25) for the plaquette edge-degree count used by
    \Cref{ch:chaining-estimate}.  Section~4, equation~(4.24)
    of that paper directly confirms the count of $2(d - 1) = 6$
    plaquettes through a given edge in dimension $d = 4$,
    underwriting $n_{\mathrm{share}} = 4$ in
    \cref{def:chain-factor}.
  \item \codepath{Seiler1982} (Lecture Notes in Physics 159):
    the standard $\muplaq = 5d - 4 = 16$ count is widely cited in
    lattice gauge theory and attributed to this monograph; the
    project's web-search pass did not surface the formula verbatim,
    so $\muplaq = 16$ is marked literature-shape (with the
    BF-conservative fallback $20\,e$ available).
\end{itemize}

The detailed audit findings sit in the project's accompanying
audit notes; the manifest above is the user-facing entry point
into that material.
